\section{METODOLOGÍA}

Para el desarrollo de la presente investigación se empleará, en una primera fase, una metodología de revisión integradora (integrative review) orientada a mapear de manera sistemática las primitivas criptográficas post-cuánticas emergentes y a evaluar su grado de adecuación para su integración con infraestructuras blockchain. Este enfoque permite sintetizar hallazgos provenientes de fuentes heterogéneas -informes de estandarización del NIST (Alagic et al., 2022; Alagic et al., 2025), estudios de evaluación de rendimiento (Lopez, Cadena y Rahman, 2025; Chaurasia, Sharma y Sharma, 2025) y propuestas de protocolos de votación post-cuántica (Ezeh et al., 2025; Yu, 2025; Wabartha, 2025)- con el fin de construir un marco de decisión fundamentado para la selección de los algoritmos y bibliotecas que conformarán el prototipo.\

En una segunda fase, el desarrollo seguirá un enfoque de investigación experimental basado en la construcción de prototipos (prototype-based experimental research), en el cual cada componente del sistema -módulo criptográfico, circuitos de conocimiento cero, contratos inteligentes e interfaz de línea de comandos- se implementa de manera incremental y se somete a pruebas unitarias antes de su integración en el pipeline completo. Este enfoque experimental resulta consistente con la naturaleza aplicada del problema de investigación, que exige no solo la formulación teórica de un protocolo, sino también la demostración empírica de su viabilidad operativa sobre hardware convencional.\

Para el módulo criptográfico post-cuántico se utilizará la biblioteca liboqs (Open Quantum Safe), reconocida como una de las implementaciones de referencia más ampliamente adoptadas en la industria y en la comunidad académica para el desarrollo con algoritmos ML-KEM y ML-DSA, accedida mediante su envoltorio (wrapper) en lenguaje Python. Para el módulo de pruebas de conocimiento cero se empleará snarkjs, herramienta que permite ejecutar el proceso de configuración confiable (Trusted Setup), generar los testigos (witnesses) a partir de las señales privadas y públicas del circuito, construir las pruebas Groth16 correspondientes y exportar automáticamente el contrato verificador en lenguaje Solidity, en combinación con Circom para la definición del circuito aritmético (Tortola et al., 2024).\

La fase de integración contempla el despliegue de los contratos inteligentes -el tablero de anuncios (Bulletin Board) y el verificador de pruebas- sobre un entorno simulado de la red Ethereum mediante Hardhat, lo cual permite validar la viabilidad técnica del sistema sin incurrir en los costos ni en la latencia asociados a una red blockchain de producción. Sobre este entorno se ejecutarán pruebas unitarias de cada componente y benchmarks de escalabilidad que simulan la concurrencia de hasta 10,000 votantes, midiendo de manera sistemática los indicadores definidos en la operacionalización de variables: tiempo de generación de claves, tiempo de generación de la prueba ZKP, costo de gas en la máquina virtual de Ethereum (EVM) y tiempo total del pipeline E2E.\

Finalmente, se realizarán pruebas unitarias de los módulos criptográficos y de conocimiento cero empleando el framework pytest, con el propósito de garantizar la corrección funcional de cada componente de manera aislada antes de proceder a las pruebas de integración end-to-end, siguiendo buenas prácticas de ingeniería de software aplicadas al desarrollo de sistemas criptográficos críticos.\

